%! Sinusoidal Functions — Unit 3, Lesson 3.5 — Warm-Up (Answer Key)
\documentclass[10pt]{article}
\usepackage{precalculus-article}
\usepackage{precalculus-key}
\newcommand{\TallMath}[1]{$\displaystyle #1\rule[-1.4em]{0pt}{3.2em}$}

\begin{document}

\pageheader{Unit 3, Lesson 3.5}{Warm-Up --- Answer Key}
\namedateperiod

\vspace{0.15in}

\begin{notesbox}{Warm-Up: Reviewing Sine and Cosine Graphs --- Key}

\textbf{1.} What is the maximum value of $f(\theta) = \sin\theta$?
The minimum value? At what $\theta$-values in $[0, 2\pi]$ does each occur?

\vspace{0.08in}
\begin{tabular}{@{}ll@{}}
  Maximum value: \ans{1} \quad occurs at $\theta = $ \ans{$\pi/2$} \\[10pt]
  Minimum value: \ans{$-1$} \quad occurs at $\theta = $ \ans{$3\pi/2$}
\end{tabular}

\vspace{0.18in}

\textbf{2.} What is the period of $g(\theta) = \cos\theta$?
How can you determine the period from looking at the graph?

\medskip
\ans{The period is $2\pi$.} The period can be determined from the graph by measuring the
horizontal distance from one point on the curve to the next point where the
curve \ans{repeats the same shape} --- for example, from one peak to the next peak,
or from one zero (with the same direction of travel) to the next such zero.

\vspace{0.18in}

\textbf{3.} True or false: $\sin(-\pi/4) = \sin(\pi/4)$. Explain your reasoning.

\medskip
\ans{False.} Because $\sin\theta$ is an odd function, $\sin(-\theta) = -\sin(\theta)$
for all $\theta$. Therefore
$\sin(-\pi/4) = -\sin(\pi/4) = $ \ans{$-\sqrt{2}/2$},
which is not equal to $\sin(\pi/4) = \sqrt{2}/2$.

\end{notesbox}

\end{document}
